\section{Introduction}
\label{sec:intro}

The rapid proliferation of connected and autonomous vehicles has positioned smart mobility as a cornerstone of Beyond 5G (B5G) and 6G cellular networks. Modern urban vehicular environments are characterized by an explosive growth in data-intensive and latency-critical applications, such as high-definition (HD) map downloading, cooperative adaptive cruise control, and augmented reality infotainment. To support these stringent Quality of Service (QoS) requirements, vehicular networks must simultaneously satisfy three highly coupled resource demands: spectrum efficiency, energy efficiency, and computation offloading.

The optimization of these three dimensions is fundamentally intertwined. For instance, a vehicle attempting to execute a computationally heavy task might choose to offload a portion of it to a Multi-Access Edge Computing (MEC) server located at a Roadside Unit (RSU). However, offloading more data requires higher uplink transmission rates, which consequently consumes more spectrum sub-bands and drains on-board battery energy. Conversely, lowering the transmission power to save energy reduces the uplink data rate, forcing the vehicle to rely on its constrained local CPU, thereby increasing computation latency. Finding a holistic equilibrium across all vehicles within this three-dimensional resource space under dynamic channel conditions and high-speed mobility remains an open challenge.

Classical mathematical techniques for resource allocation, such as convex optimization and Lyapunov drift frameworks, typically require perfect network-wide Channel State Information (CSI) and suffer from severe computational overhead. As vehicular topologies change by the millisecond, centralized iterative solvers fail to compute actionable policies within the coherence time of the vehicular channels.

To overcome the intractability of mathematical programming, deep reinforcement learning (DRL) algorithms have been widely adopted. However, single-agent architectures scale poorly in dense networks. Recent advances in Multi-Agent Reinforcement Learning (MARL), including MAPPO and MADDPG, have digitized distributed vehicular decisions but suffer from flat state representations. Such architectures struggle to capture the complex, dynamic interference graphs formed by moving vehicles, leading to models that memorize specific spatial layouts but completely fail to generalize across arbitrary urban topologies. Furthermore, algorithms utilizing continuous action spaces often employ deterministic gradients that restrict wide state-space exploration, yielding suboptimal convergence.

To directly address these critical shortcomings, we propose a novel Hierarchical Graph Attention Multi-Agent Soft Actor-Critic (HiGAT-MASAC) framework to orchestrate joint spectrum, power, and edge computing allocation in Next-Generation Vehicular Networks. Our design disentangles the macroscopic cellular bandwidth constraints from the rapid topological interference shifts through the following key innovations:

\begin{itemize}
    \item \textbf{Hierarchical Timescale Decomposition:} The complex resource optimization problem is structured as a two-tier Decentralized Partially Observable Markov Decision Process (Dec-POMDP). RSUs operate at a slow macroscopic timescale (e.g., 100 ms) to allocate spectrum sub-bands and MEC processing capacity, while individual vehicles operate at a fast microscopic timescale (e.g., 10 ms) to optimize continuous transmit power and computation offload ratios within their granted resource envelopes.
    \item \textbf{Dynamic Graph Attention Networks (GAT):} We model the transient Vehicle-to-Vehicle (V2V) interference relationships as dynamic graphs. A Multi-Head GAT layer automatically learns attention weights, intelligently emphasizing high-interference neighbors over distant anomalies without relying on manual feature engineering. This inherently enables the policy to generalize seamlessly across unseen vehicular clustering configurations.
    \item \textbf{Hybrid Maximum-Entropy RL Optimization:} By integrating the Soft Actor-Critic (SAC) framework across both tiers, the architecture robustly handles hybrid action parameters. The continuous SAC networks are augmented with Gumbel-Softmax discrete relaxations for the RSU tier, whilst maximizing policy entropy ensures robust, exploratory gradient flows that surpass traditional deterministic off-policy architectures.
\end{itemize}

Specifically, the major contributions of this paper are summarized as follows:
\begin{enumerate}
    \item We mathematically formulate the joint resource optimization challenge---encompassing system spectrum mapping, task delay limits, and offloading energy bounds---as an NP-hard mixed-integer non-linear program.
    \item We design HiGAT-MASAC, a pioneering two-tier hierarchical MARL algorithm that physically segregates infrastructural allocations from distributed vehicle executions to effectively cut transmission overhead and improve control scalability.
    \item We organically integrate Graph Attention layers into the SAC value functions. The localized topological representations passed across the tiers bypass the fundamental bottleneck of purely fully connected neural networks seen in contemporary MARL IoV systems.
    \item We deploy an extensive suite of empirical Manhattan grid simulations compliant with 3GPP UMi mobility channels. Evaluation results emphatically validate that HiGAT-MASAC guarantees superior delay-reduction, maximal sum-rate throughputs, and highly favorable energy reductions universally compared against state-of-the-art baseline models including GNN-DDQN, MAPPO, and MADDPG.
\end{enumerate}

The remainder of this paper is structured as follows. Section \ref{sec:related} reviews pertinent literature across multi-agent vehicular learning and graph neural communications. Section \ref{sec:model} formulates the comprehensive network, channel, and latency system models utilized. Section \ref{sec:algo} details the complete derivation of the proposed HiGAT-MASAC coordination framework. Section \ref{sec:eval} illustrates the extensive quantitative results and analysis. Finally, Section \ref{sec:conc} concludes the paper with directions for future exploration.
